\chapter{Antecedentes}

Este capítulo ofrece un panorama general de herramientas de generación automática de código orientadas a hardware reconfigurable. Se destacan los enfoques de \textit{High-Level Synthesis} (HLS) para traducir descripciones de alto nivel a HDL, su presencia en la industria y sus principales aplicaciones.

\section{Síntesis de alto nivel y herramientas de generación de código en la industria}

La síntesis de alto nivel (HLS) busca reducir la brecha entre la especificación algorítmica y la implementación física en FPGA, generando descripciones RTL (VHDL o Verilog) a partir de modelos con mayor nivel de abstracción. En este enfoque, la herramienta se encarga de tareas estructurales como planificación temporal, inserción de \textit{pipeline}, reparto de operadores y ajuste de paralelismo. El diseñador, por su parte, mantiene el control de decisiones de arquitectura: precisión numérica, latencia objetivo, estrategia de memoria y límites de iteración \cite{coussy2009high,gaide2011high}.

En plataformas reconfigurables actuales, HLS se utiliza tanto para prototipado rápido como para acelerar etapas de diseño donde la codificación RTL manual sería costosa en tiempo. Su valor práctico no está solo en ``generar HDL'', sino en permitir ciclos iterativos de exploración: cambiar una restricción, regenerar, comparar recursos y validar nuevamente la funcionalidad. Esa dinámica es particularmente útil cuando se debe balancear desempeño y área en diseños con requerimientos estrictos de temporización.

En este contexto, el flujo moderno de desarrollo no termina en la síntesis. El valor real se obtiene al cerrar el ciclo con verificación funcional y temporal sobre hardware, por ejemplo mediante co-simulación y \textit{FPGA-in-the-Loop}. Estas etapas permiten detectar inconsistencias que no son evidentes en simulación de alto nivel, como errores por cuantización, latencias desalineadas o efectos de saturación acumulada en lazo cerrado.

Finalmente, HLS también ha cambiado la forma de documentar diseños: además del código generado, se reportan métricas de trazabilidad (configuración usada, reportes de recursos, latencias, restricciones de reloj y casos de prueba), lo que facilita reproducibilidad y comparación entre variantes de un mismo algoritmo.

\section{Herramientas de MATLAB para generación y verificación de código}

El ecosistema de MathWorks para FPGA se apoya principalmente en tres componentes complementarios: modelado en MATLAB/Simulink como referencia funcional, ajuste numérico con \textit{fixed-point}, y generación/verificación de HDL mediante \textit{HDL Coder} y \textit{HDL Verifier}. Este enfoque permite recorrer un flujo continuo desde el diseño algorítmico hasta la validación sobre hardware real, manteniendo una única fuente de verdad para comparar resultados.

Una ventaja concreta de este esquema es la trazabilidad entre decisiones de modelado y consecuencias físicas. Por ejemplo, modificar el número de bits fraccionarios o activar compartición de recursos no solo altera una simulación: también cambia latencia, utilización de operadores y exigencia de reloj interno en la implementación final. Esta relación directa entre modelo y reporte de síntesis es uno de los elementos más útiles para iterar con criterio de ingeniería.

\subsection{Fixed-point designer}

Antes de generar HDL, el modelo debe cumplir condiciones de sintetizabilidad: tamaños estáticos, lazos acotados, rutas de control deterministas y tipos de dato compatibles con hardware. En la práctica, esta etapa es donde más tiempo se invierte, porque define la robustez del flujo completo. Un modelo que ``simula bien'' en alto nivel puede fallar en síntesis si mantiene dinámicas ambiguas, conversiones implícitas o rangos numéricos mal acotados.

Para evitar esas fallas, se recomienda fijar explícitamente escalas, redondeo y saturación en bloques críticos, y validar la diferencia entre referencia y modelo cuantizado con bancos de prueba representativos. En este proyecto, esta práctica fue clave en ejemplos de control y optimización, donde pequeñas desviaciones numéricas pueden modificar de forma importante la respuesta en lazo cerrado.

\subsection{HDL Coder}

HDL Coder permite generar RTL (VHDL o Verilog) desde modelos de Simulink o funciones de MATLAB. El flujo se apoya en el \textit{HDL Workflow Advisor}, que configura el tipo de objetivo, los relojes, el \textit{oversampling} y las optimizaciones de arquitectura. Entre sus capacidades destacadas se incluyen: generación automática de \textit{pipelines}, reparto de recursos en bucles (\textit{resource sharing}), inserción de registros para cumplir temporización, y reportes de latencia y utilización. La herramienta exige que los modelos sean sintetizables: dimensiones estáticas, ciclos acotados y tipos de datos compatibles (preferentemente \textit{fixed-point}) \cite{matlabHDLCoder}.

En la práctica, HDL Coder permite cerrar un ciclo iterativo de modelado, generación y revisión de resultados de síntesis. Este ciclo resulta más efectivo cuando se mantiene una disciplina de configuración estable entre iteraciones, de modo que las diferencias observadas puedan atribuirse a cambios de arquitectura y no a cambios del procedimiento.

Un elemento crítico es la conversión de punto flotante a punto fijo. HDL Coder se integra con Fixed-Point Designer para proponer tamaños de palabra, evaluar error numérico y definir reglas de saturación y redondeo. Esta etapa determina el rango dinámico y la precisión, afectando directamente el uso de recursos y la equivalencia funcional. En consecuencia, la recomendación general es validar primero el modelo en punto fijo y luego generar HDL, evitando conversiones tardías que puedan introducir saturaciones no controladas \cite{fixedPointDesigner}.

Además, HDL Coder puede generar IP cores con interfaces estándar para integrarlos en flujos de diseño de FPGA. Este paso es relevante cuando el diseño se conecta a buses o a procesadores embebidos, ya que permite empaquetar el bloque como IP con interfaces configurables y documentación de puertos.

\subsection{HDL Verifier}

HDL Verifier complementa el flujo de HDL Coder con mecanismos de verificación. Permite realizar co-simulación entre Simulink y simuladores HDL de terceros, comparando la salida del modelo con la del HDL generado y reportando discrepancias. Esta etapa es especialmente útil para detectar errores de cuantización, latencias desalineadas o efectos de saturación que no aparecen en simulación puramente en software \cite{matlabHDLVerifier}.

Una característica central es \textit{FPGA-in-the-Loop} (FIL), donde el diseño HDL se sintetiza, se implementa en FPGA y se ejecuta en hardware real mientras Simulink actúa como banco de pruebas. El modo \textit{lockstep} sincroniza ciclos de FPGA con pasos de simulación; el modo \textit{free-running} permite que el hardware ejecute a su reloj interno y Simulink intercambie datos bajo demanda. El ajuste del \textit{overclocking factor} es esencial para alinear el número de ciclos internos con el muestreo de Simulink, particularmente cuando el HDL tiene latencias multicíclo.

Otro parámetro relevante en FIL es el \textit{output delay}. Este valor indica cuántos periodos de muestreo de entrada deben transcurrir antes de obtener una salida válida, independiente del reloj interno. En un sistema retroalimentado con tiempo de muestreo fijo, dicho retardo se traduce en una demora efectiva dentro del lazo. Por ejemplo, si el controlador fue diseñado para leer la planta cada 1 ms y actuar en ese mismo periodo, un \textit{output delay} de 4 implica que el actuador solo recibe una acción cada 4 ms; la planta cambia cada 1 ms, pero el control solo se actualiza cada cuatro muestras. Esto equivale a introducir un retardo de 4 ms en el lazo, lo que reduce el margen de fase y puede degradar el desempeño o incluso inestabilizar el sistema. En esos casos, se debe reconsiderar la arquitectura (reducir latencia, ajustar el \textit{overclocking}, o rediseñar el controlador para el nuevo periodo efectivo) antes de validar el comportamiento en hardware.

\subsection{Discusión}

Dentro del alcance de esta memoria, HDL Coder y HDL Verifier conforman un flujo cerrado: generación RTL, análisis de arquitectura, verificación funcional y validación temporal en FPGA. La experiencia práctica de los ejemplos confirma que la principal fuente de éxito o fracaso no es la herramienta en sí, sino la disciplina de modelado aplicada desde el inicio.

En términos metodológicos, se identifican cuatro criterios que se repiten en todos los casos exitosos: (i) definición temprana de tipos y escalas, (ii) control explícito de latencia mediante \textit{pipeline}/\textit{resource sharing}, (iii) coherencia entre muestreo del modelo y temporalidad interna del HDL, y (iv) validación incremental con bancos de prueba trazables. Cuando alguno de estos criterios se omite, la probabilidad de retrabajo aumenta de manera significativa.

Por esta razón, el aporte de este capítulo no es solo describir herramientas, sino establecer una base de diseño reproducible orientada a HDL, sobre la cual se construyen los resultados experimentales de los capítulos siguientes.
